% Generated draft: requires mathematical and source review.
\documentclass[11pt]{article}
\usepackage[margin=25mm]{geometry}
\usepackage{fontspec}
\setmainfont{Latin Modern Roman}
\usepackage{amsmath,amssymb,mathtools,unicode-math}
\setmathfont{Latin Modern Math}
\usepackage{enumitem,needspace,xurl,hyperref}
\setlength{\parindent}{0pt}
\setlength{\parskip}{5pt}
\setlength{\emergencystretch}{3em}
\newcommand{\N}{\mathbb{N}}
\newcommand{\Z}{\mathbb{Z}}
\newcommand{\Q}{\mathbb{Q}}
\newcommand{\R}{\mathbb{R}}
\newcommand{\C}{\mathbb{C}}
\newcommand{\F}{\mathbb{F}}
\newcommand{\A}{\mathbb{A}}
\newcommand{\PP}{\mathbb{P}}
\newcommand{\E}{\mathbb{E}}
\newcommand{\eps}{\varepsilon}
\newcommand{\dd}{\,\mathrm{d}}
\begin{document}
\section*{UnsolvedMath ID: 800866}
\textbf{Source problem code:} catalogue entry
\subsection{Definitions and mathematical statement}
Use the conventional real interpretation of a positive definite symmetric matrix: a matrix \(A=(a_{ij})\in\R^{5\times5}\) satisfying \(A^{\mathsf T}=A\) and \(x^{\mathsf T}Ax>0\) for every \(x\in\R^5\setminus\{0\}\). The supplied statement does not explicitly specify the scalar field.

Write \([5]=\{1,2,3,4,5\}\). For equal-sized subsets \(I,J\subseteq[5]\), let \(A_{I,J}\) denote the submatrix with rows indexed by \(I\) and columns indexed by \(J\), both in increasing order. An almost-principal minor is \(\det A_{I,J}\) with \(|I\setminus J|=|J\setminus I|=1\). Thus it can be indexed by distinct indices \(i,j\) and a set \(K\subseteq[5]\setminus\{i,j\}\), with \(I=K\cup\{i\}\) and \(J=K\cup\{j\}\). These minors have sizes \(1\) through \(4\). Since \(A\) is symmetric, interchanging \(I\) and \(J\) gives the same determinant, so use the index set
\[
\mathcal T=\{(i,j,K):1\le i<j\le5,\ K\subseteq[5]\setminus\{i,j\}\}.
\]
For \(t=(i,j,K)\in\mathcal T\), put
\[
m_t(A)=\det A_{K\cup\{i\},K\cup\{j\}}.
\]

Under the exact-vanishing-pattern interpretation, the problem attributed to Sturmfels (2007) is to characterize all subsets \(S\subseteq\mathcal T\) for which there exists a real symmetric positive definite \(5\times5\) matrix \(A\) such that
\[
\forall t\in\mathcal T,\qquad m_t(A)=0\ \Longleftrightarrow\ t\in S.
\]
The source asks which sets of almost-principal minors ``can be zero'' without explicitly requiring the remaining minors to be nonzero. If only simultaneous vanishing is intended, the target instead is to characterize those \(S\subseteq\mathcal T\) for which such a matrix exists satisfying \(m_t(A)=0\) for every \(t\in S\), with no condition on \(t\notin S\). This distinction remains unresolved by the supplied record.
\subsection{Short English statement}
Determine which patterns of vanishing almost-principal minors occur in positive definite symmetric \(5\times5\) matrices, using the conventional real interpretation. The source leaves unclear whether the specified set must be the entire zero pattern or may allow additional vanishing minors.
\subsection{Sources}
\begin{itemize}
\item Local UnsolvedMath registry: permanent ID \(800866\), exact source problem code \texttt{catalogue entry}, title \emph{Open Problems in Algebraic Statistics --- Specific Problem 5}. The supplied record attributes the problem to Sturmfels, with proposed year 2007, and labels its status as open. This label is supplied registry information, not an independently verified mathematical or literature assessment.
\item Supplied source: Sturmfels, \emph{Open Problems in Algebraic Statistics} (2007), Specific Problem \#5, \url{https://arxiv.org/abs/0707.4558}. The provenance records access on 2026-10-16 and extraction method \texttt{arxiv\_latex}; these are supplied provenance details, not a claim of independent access or verification.
\end{itemize}
\end{document}
